\documentclass{article}
\usepackage{amsmath,amssymb,amsthm}

\newtheorem{definition}{Definition}
\newcommand{\source}[1]{\par\noindent\textbf{Source.} #1}
\newcommand{\lean}[1]{\par\noindent\textbf{Lean target.} \texttt{#1}}
\newcommand{\status}[1]{\par\noindent\textbf{Status.} \texttt{\detokenize{#1}}}

\begin{document}

\section*{Blueprint 004: Represented Region von Neumann Algebra}

\paragraph{Superseded note.}
This historical draft is superseded by
\texttt{blueprints/010-gns-represented-region-layer.tex}.  It is preserved
only as provenance for the earlier draft and should not be used as an active
Aristotle source.

\begin{definition}[RegionVNPackage, superseded ID: VN-001]
Let \((H_\omega,\pi_\omega,\Omega_\omega)\) be the GNS representation of a
state \(\omega\) on a quasi-local algebra \(A\).  For a region \(\Gamma\), let
\(A_\Gamma\subset A\) be the region subalgebra.  The represented region von
Neumann algebra is
\[
  M_\Gamma(\omega) := \pi_\omega(A_\Gamma)'' \subset B(H_\omega),
\]
the bicommutant of the represented region algebra.
\source{Concrete von Neumann algebra by bicommutant; mathlib
\texttt{Mathlib.Analysis.VonNeumannAlgebra.Basic}, especially
\texttt{VonNeumannAlgebra}, \texttt{VonNeumannAlgebra.commutant}, and
\texttt{VonNeumannAlgebra.commutant\_commutant}.  Project proof note
\texttt{runs/2026-06-18-stabilizer-algebra-classification/artifacts/004-internal-stabilizer-reduction-type-classification-R1.tex},
Theorem 5.1, uses this region algebra.}
\lean{Formalization.OperatorAlgebra.RegionVNPackage}
\status{blueprint_draft_source_locator_fixed}
\end{definition}

\end{document}
